% ======================================================================
%  Taller 1 -- Punto 1.1 (cuerpo)
%  Requiere el bloque "preambulo_taller.tex" en el preámbulo.
%  Se puede pegar directamente en el documento o incluir con
%  \input{punto_1_1_cuerpo}.
% ======================================================================
% ====================== INICIO DEL PUNTO 1.1 ===========================



%======================================================================
\subsection*{El conjunto $V$}
%======================================================================

\begin{definicionT}[Combinación lineal y conjunto generado]
Dadas funciones $\varphi_1,\dots,\varphi_k$ definidas en $[0,2\pi]$, una
\emph{combinación lineal con coeficientes reales} de ellas es una función de la forma
$c_1\varphi_1+c_2\varphi_2+\dots+c_k\varphi_k$, con $c_1,\dots,c_k\in\mathbb R$. El conjunto de
todas estas combinaciones se llama \emph{conjunto generado} por $\varphi_1,\dots,\varphi_k$ y
se denota $\spanR\{\varphi_1,\dots,\varphi_k\}$.
\end{definicionT}

Según esta definición, los elementos de $V$ son las funciones de la forma
\begin{equation}\label{eq:forma-general}
  f(t) = a_1\Cu + b_1\Su + a_2\Cd + b_2\Sd,
  \qquad a_1, b_1, a_2, b_2 \in \mathbb{R}.
\end{equation}
En adelante, $a_n$ denota el coeficiente de $\cos nt$ y $b_n$ el de $\sin nt$. El subíndice
$n$ es la \emph{frecuencia} de la función, es decir, el número de oscilaciones completas que
realiza en el intervalo $[0,2\pi]$; así, $\Cu$ y $\Su$ tienen frecuencia $1$, y $\Cd$ y $\Sd$
tienen frecuencia $2$ (figura~\ref{fig:generadoras}). Todo elemento de $V$ es una función que
toma valores reales.

\begin{figure}[H]
\centering
\begin{tikzpicture}
\begin{axis}[ondaT, width=14cm, height=5.8cm,
  xtick={0,pi/4,pi/2,pi,3*pi/2,2*pi},
  xticklabels={$0$,$\frac{\pi}{4}$,$\frac{\pi}{2}$,$\pi$,$\frac{3\pi}{2}$,$2\pi$},
  legend style={at={(0.5,1.04)}, anchor=south, legend columns=4, draw=none,
                font=\small, column sep=8pt},
  ]
  \addplot[very thick, colCu] {cos(x)};    \addlegendentry{$\cos t$}
  \addplot[very thick, colSu] {sin(x)};    \addlegendentry{$\sin t$}
  \addplot[very thick, colCd] {cos(2*x)};  \addlegendentry{$\cos 2t$}
  \addplot[very thick, colSd] {sin(2*x)};  \addlegendentry{$\sin 2t$}
  \draw[dashed, black!50] (axis cs:0.7854,-1.25) -- (axis cs:0.7854,1.25);
  \draw[dashed, black!50] (axis cs:1.5708,-1.25) -- (axis cs:1.5708,1.25);
  \draw[dashed, black!50] (axis cs:3.1416,-1.25) -- (axis cs:3.1416,1.25);
\end{axis}
\end{tikzpicture}
\caption{Las cuatro funciones que generan $V$. Las de frecuencia $1$ completan una oscilación
en $[0,2\pi]$ y las de frecuencia $2$ completan dos. Las líneas punteadas marcan los puntos
$t=\pi/4,\ \pi/2,\ \pi$ que, junto con $t=0$, se usan en la demostración del
lema~\ref{lema:independencia}.}
\label{fig:generadoras}
\end{figure}

\begin{definicionT}[Independencia lineal sobre $\mathbb R$]
Las funciones $\varphi_1,\dots,\varphi_k$ son \emph{linealmente independientes sobre
$\mathbb R$} si la única combinación lineal con coeficientes reales que da la función cero (la
que vale $0$ en todo $t$) es aquella con todos sus coeficientes iguales a cero; es decir, si
\[
  c_1\varphi_1(t)+\dots+c_k\varphi_k(t)=0 \ \text{ para todo } t\in[0,2\pi]
  \quad\Longrightarrow\quad c_1=\dots=c_k=0 .
\]
\end{definicionT}

\begin{lemaT}{lema:independencia}
Las funciones $\cos t$, $\sin t$, $\cos 2t$ y $\sin 2t$ son linealmente independientes sobre
$\mathbb R$. En consecuencia, si dos expresiones de la forma $f(t)=a_1\cos t+b_1\sin t+a_2\cos 2t+b_2\sin 2t$ \eqref{eq:forma-general}
representan la misma función, sus coeficientes coinciden.
\end{lemaT}

\begin{proof}
Supóngase que
\begin{equation}\label{eq:combinacion-cero}
  a_1\cos t + b_1\sin t + a_2\cos 2t + b_2\sin 2t = 0
  \qquad\text{para todo } t\in[0,2\pi].
\end{equation}
Como la igualdad se cumple para todo $t$, se cumple en particular en los puntos donde varias
de las cuatro funciones valen $0$ o $\pm 1$, que son los marcados en la
figura~\ref{fig:generadoras}. Los valores en esos puntos son
\[
\begin{array}{c|cccc}
 t & \cos t & \sin t & \cos 2t & \sin 2t\\ \hline
 0 & 1 & 0 & 1 & 0\\
 \pi/4 & \tfrac{\sqrt2}{2} & \tfrac{\sqrt2}{2} & 0 & 1\\
 \pi/2 & 0 & 1 & -1 & 0\\
 \pi & -1 & 0 & 1 & 0
\end{array}
\]
Al evaluar $a_1\cos t+b_1\sin t+a_2\cos 2t+b_2\sin 2t=0$ \eqref{eq:combinacion-cero} en $t=0$ y en $t=\pi$ se obtienen las ecuaciones
$a_1+a_2=0$ y $-a_1+a_2=0$. Sumándolas resulta $2a_2=0$, es decir, $a_2=0$, y entonces
$a_1=0$. Al evaluar en $t=\pi/2$ se obtiene $b_1-a_2=0$, y como $a_2=0$, también $b_1=0$.
Finalmente, al evaluar en $t=\pi/4$ queda
$\tfrac{\sqrt2}{2}a_1+\tfrac{\sqrt2}{2}b_1+b_2=0$, y como $a_1=b_1=0$, resulta $b_2=0$. Así,
los cuatro coeficientes son cero.

Para la segunda afirmación, supóngase que una misma función $f$ se escribe con coeficientes
$a_1,b_1,a_2,b_2$ y también con coeficientes $\tilde a_1,\tilde b_1,\tilde a_2,\tilde b_2$. Al
restar ambas expresiones se obtiene
$(a_1-\tilde a_1)\cos t+(b_1-\tilde b_1)\sin t+(a_2-\tilde a_2)\cos 2t+(b_2-\tilde b_2)\sin 2t=0$
para todo $t$, y por la primera parte cada diferencia es cero, es decir, $a_1=\tilde a_1$,
$b_1=\tilde b_1$, $a_2=\tilde a_2$ y $b_2=\tilde b_2$.
\end{proof}

Gracias al lema tiene sentido hablar de \emph{los} coeficientes de un elemento de $V$, y dos
elementos de $V$ son iguales exactamente cuando sus cuatro coeficientes coinciden. Este es el
criterio que se usa en todas las verificaciones siguientes.

%======================================================================
\subsection*{Propiedades de espacio vectorial complejo}
%======================================================================

\begin{definicionT}[Espacio vectorial complejo]
Un conjunto $V$ es un \emph{espacio vectorial complejo} si tiene definidas dos operaciones, la
suma $f+g$ de dos elementos y la multiplicación $z\cdot f$ de un elemento por un número
complejo $z$, cuyos resultados pertenecen otra vez a $V$ y que cumplen las
siguientes propiedades para todo $f, g, h\in V$ y todo $z, w\in\mathbb{C}$.

\smallskip
\noindent\textbf{Propiedades de la suma.}
\begin{enumerate}
  \item[(S1)] $f+g = g+f$.
  \item[(S2)] $(f+g)+h = f+(g+h)$.
  \item[(S3)] Existe un elemento $0\in V$ tal que $f+0 = f$.
  \item[(S4)] Para cada $f$ existe $-f\in V$ tal que $f+(-f) = 0$.
\end{enumerate}
\noindent\textbf{Propiedades de la multiplicación por escalares.}
\begin{enumerate}
  \item[(M1)] $z\cdot(f+g) = z\cdot f + z\cdot g$.
  \item[(M2)] $(z+w)\cdot f = z\cdot f + w\cdot f$.
  \item[(M3)] $(zw)\cdot f = z\cdot(w\cdot f)$.
  \item[(M4)] $1\cdot f = f$.
\end{enumerate}
\end{definicionT}

En $V$ la suma es la suma usual de funciones, $(f+g)(t)=f(t)+g(t)$, mientras que la
multiplicación por escalares complejos se construye a partir de la regla del enunciado. Las dos
operaciones se estudian por separado.

\subsubsection*{La suma}

Sean $f, g, h\in V$. Por la definición de $V$, existen números reales
$a_1, b_1, a_2, b_2$, $a_1', b_1', a_2', b_2'$ y $a_1'', b_1'', a_2'', b_2''$ tales que
\begin{align*}
  f &= a_1\cos t + b_1\sin t + a_2\cos 2t + b_2\sin 2t,\\
  g &= a_1'\cos t + b_1'\sin t + a_2'\cos 2t + b_2'\sin 2t,\\
  h &= a_1''\cos t + b_1''\sin t + a_2''\cos 2t + b_2''\sin 2t.
\end{align*}

\paragraph{La suma pertenece a $V$.}
Al sumar $f$ y $g$ y agrupar los términos que contienen la misma función, se obtiene
\begin{equation}\label{eq:suma}
  f+g = (a_1+a_1')\cos t + (b_1+b_1')\sin t + (a_2+a_2')\cos 2t + (b_2+b_2')\sin 2t.
\end{equation}
Como $a_1+a_1'$, $b_1+b_1'$, $a_2+a_2'$ y $b_2+b_2'$ son números reales, $f+g$ es de la
forma $f(t)=a_1\cos t+b_1\sin t+a_2\cos 2t+b_2\sin 2t$ \eqref{eq:forma-general} y por tanto pertenece a $V$. En otras palabras, sumar dos
elementos de $V$ equivale a sumar sus coeficientes uno a uno.

\paragraph{(S1) $f+g = g+f$.}
Por $f+g=(a_1+a_1')\cos t+(b_1+b_1')\sin t+(a_2+a_2')\cos 2t+(b_2+b_2')\sin 2t$ \eqref{eq:suma}, los coeficientes de $f+g$ son $a_1+a_1'$, $b_1+b_1'$, $a_2+a_2'$ y
$b_2+b_2'$, mientras que los de $g+f$ son $a_1'+a_1$, $b_1'+b_1$, $a_2'+a_2$ y $b_2'+b_2$. Como
la suma de números reales no depende del orden, $a_1+a_1'=a_1'+a_1$, y lo mismo ocurre con los
otros tres coeficientes. Al tener los mismos coeficientes, $f+g=g+f$.

\paragraph{(S2) $(f+g)+h = f+(g+h)$.}
Al aplicar $f+g=(a_1+a_1')\cos t+(b_1+b_1')\sin t+(a_2+a_2')\cos 2t+(b_2+b_2')\sin 2t$ \eqref{eq:suma} dos veces, el coeficiente de $\cos t$ en $(f+g)+h$ es
$(a_1+a_1')+a_1''$ y en $f+(g+h)$ es $a_1+(a_1'+a_1'')$. Estos números son iguales porque la
suma de números reales es asociativa. El mismo argumento se aplica a los coeficientes de
$\sin t$, $\cos 2t$ y $\sin 2t$, de modo que $(f+g)+h = f+(g+h)$.

\paragraph{(S3) Existe $0\in V$ tal que $f+0 = f$.}
Sea $0$ la función con sus cuatro coeficientes iguales a cero,
\[
  0 = 0\cos t + 0\sin t + 0\cos 2t + 0\sin 2t,
\]
que pertenece a $V$ porque $0$ es un número real; se trata de la función que vale $0$ en todo
$t$. Por $f+g=(a_1+a_1')\cos t+(b_1+b_1')\sin t+(a_2+a_2')\cos 2t+(b_2+b_2')\sin 2t$ \eqref{eq:suma},
\[
  f+0 = (a_1+0)\cos t + (b_1+0)\sin t + (a_2+0)\cos 2t + (b_2+0)\sin 2t = f.
\]

\paragraph{(S4) Para cada $f$ existe $-f\in V$ tal que $f+(-f) = 0$.}
Sea $-f$ la función con los coeficientes de $f$ cambiados de signo,
\[
  -f = (-a_1)\cos t + (-b_1)\sin t + (-a_2)\cos 2t + (-b_2)\sin 2t,
\]
que pertenece a $V$ porque $-a_1, -b_1, -a_2, -b_2$ son números reales. Por $f+g=(a_1+a_1')\cos t+(b_1+b_1')\sin t+(a_2+a_2')\cos 2t+(b_2+b_2')\sin 2t$ \eqref{eq:suma},
\[
  f+(-f) = (a_1-a_1)\cos t + (b_1-b_1)\sin t + (a_2-a_2)\cos 2t + (b_2-b_2)\sin 2t = 0.
\]
Con esto, la suma en $V$ produce elementos de $V$ y cumple (S1)--(S4).

\subsubsection*{La multiplicación por escalares complejos}

La regla del enunciado está escrita en términos de la forma polar de un número complejo, por
lo que conviene fijar primero la notación correspondiente.

\begin{definicionT}[Partes, conjugado y módulo de un número complejo]
Todo número complejo se escribe como $z=a+ib$, con $a,b\in\mathbb R$ e $i^2=-1$. El número
$a=\Real z$ es su \emph{parte real} y $b=\Imag z$ su \emph{parte imaginaria}. El
\emph{conjugado} de $z$ es $\overline z=a-ib$ y su \emph{módulo} es $|z|=\sqrt{a^2+b^2}$. Al
multiplicar $z$ por su conjugado se obtiene
\begin{equation}\label{eq:modulo}
  z\,\overline z=(a+ib)(a-ib)=a^2-i^2b^2=a^2+b^2=|z|^2 .
\end{equation}
El conjugado respeta la suma y el producto y, al aplicarse dos veces, devuelve el número
original: para $z,w\in\mathbb C$,
\begin{equation}\label{eq:conjugado}
  \overline{z+w}=\overline z+\overline w,\qquad
  \overline{zw}=\overline z\,\overline w,\qquad
  \overline{\overline z}=z .
\end{equation}
Estas igualdades se comprueban escribiendo $z=a+ib$, $w=c+id$ y desarrollando ambos lados.
\end{definicionT}

\begin{definicionT}[Forma polar]
La \emph{fórmula de Euler} establece que, para $x\in\mathbb R$,
$e^{ix}=\cos x+i\sin x$. Todo número complejo $z=a+ib$ puede escribirse en \emph{forma polar}
$z=Ae^{ix}$, donde $A=|z|\ge 0$ es su módulo y $x$ es el ángulo que forma el segmento que va
de $0$ a $z$ con el eje real. Leyendo la figura como un triángulo
rectángulo de hipotenusa $A$, se tiene
\begin{equation}\label{eq:polar}
  a=A\cos x,\qquad b=A\sin x .
\end{equation}
El ángulo $x$ no es único: $x$ y $x+2\pi$ describen el mismo número, y si $z=0$ cualquier
ángulo sirve.
\end{definicionT}


\paragraph{Significado de la regla.}
La regla del enunciado,
\[
  (Ae^{ix})\cdot\cos nt = A\cos(nt-x),\qquad n=1,2,
\]
tiene una lectura gráfica directa. Como $A\cos(nt-x)=A\cos\bigl(n(t-x/n)\bigr)$, la función
resultante es la gráfica de $\cos nt$ desplazada $x/n$ unidades hacia la derecha y estirada
verticalmente por el factor $A$. El módulo del número complejo controla la altura de la onda y
el ángulo controla su posición.

\paragraph{La regla en términos de las partes real e imaginaria.}
Sea $z=a+ib\in\mathbb C$ escrito en forma polar como $z=Ae^{ix}$. Al aplicar la regla, usar la
fórmula del coseno de una resta, $\cos(nt-x)=\cos nt\cos x+\sin nt\sin x$, y luego
$a=A\cos x,\quad\penalty0 b=A\sin x$ \eqref{eq:polar}, se obtiene
\begin{align*}
  (a+ib)\cdot\cos nt
  &=A\cos(nt-x)
   =A\bigl(\cos nt\cos x+\sin nt\sin x\bigr)\\
  &=(A\cos x)\cos nt+(A\sin x)\sin nt
   =a\cos nt+b\sin nt .
\end{align*}
Por tanto,
\begin{equation}\label{eq:regla-coef}
  \destacado{(a+ib)\cdot\cos nt=a\cos nt+b\sin nt,\qquad n=1,2.}
\end{equation}
El lado derecho depende únicamente de $a$ y $b$, no del ángulo $x$ elegido, así que la regla
asigna una sola función a cada número complejo aunque su forma polar no sea única. Además, si
$z=a$ es real ($b=0$), la regla da $a\cdot\cos nt=a\cos nt$, que es el producto usual de un
número real por una función; la nueva multiplicación extiende, por tanto, la multiplicación
por números reales. En el caso particular $z=i$, es decir, $a=0$ y $b=1$, la fórmula
$(a+ib)\cdot\cos nt=a\cos nt+b\sin nt$ \eqref{eq:regla-coef} da
\begin{equation}\label{eq:i-cos}
  i\cdot\cos nt=\sin nt ,
\end{equation}
que es lo que muestra en la figura~\ref{fig:regla}.

\begin{figure}[H]
\centering
\begin{tikzpicture}
\begin{axis}[ondaT, title={$i\cdot\cos t=\cos(t-\pi/2)=\sin t$}]
  \addplot[very thick, colCu] {cos(x)};
  \addplot[very thick, colSu] {sin(x)};
  \draw[-{Stealth[length=2mm]}, thick] (axis cs:0,1.2) -- (axis cs:1.5708,1.2)
     node[midway, above, font=\footnotesize] {$\pi/2$};
  \node[colCu, font=\footnotesize] at (axis cs:5.75,1.15) {$\cos t$};
  \node[colSu, font=\footnotesize] at (axis cs:4.0,-1.2) {$\sin t$};
\end{axis}
\end{tikzpicture}
\caption{}
\label{fig:regla}
\end{figure}

\paragraph{Coordenadas complejas de un elemento de $V$.}
La regla solo indica cómo multiplicar $\cos t$ y $\cos 2t$ por un número complejo; para
multiplicar un elemento cualquiera de $V$ hace falta extenderla. La fórmula
$(a+ib)\cdot\cos nt=a\cos nt+b\sin nt$ \eqref{eq:regla-coef}, leída de derecha a izquierda, muestra cómo hacerlo, pues dice que una
pareja de términos con la misma frecuencia es un múltiplo complejo de $\cos nt$:
\[
  a_1\cos t+b_1\sin t=(a_1+ib_1)\cdot\cos t,
  \qquad
  a_2\cos 2t+b_2\sin 2t=(a_2+ib_2)\cdot\cos 2t .
\]

\begin{definicionT}[Coordenadas complejas]
Para $f\in V$ con coeficientes $a_1,b_1,a_2,b_2$ como en $f(t)=a_1\cos t+b_1\sin t+a_2\cos 2t+b_2\sin 2t$ \eqref{eq:forma-general}, sus
\emph{coordenadas complejas} son los números
\[
  \alpha=a_1+ib_1,\qquad \beta=a_2+ib_2 .
\]
La parte real de $\alpha$ es el coeficiente de $\cos t$ y su parte imaginaria es el
coeficiente de $\sin t$; lo mismo ocurre con $\beta$ para la frecuencia $2$.
\end{definicionT}

Con esta notación, todo elemento de $V$ se escribe como
\begin{equation}\label{eq:alfa-beta}
  f=\alpha\cdot\Cu+\beta\cdot\Cd .
\end{equation}
Por el lema~\ref{lema:independencia}, los números $\alpha$ y $\beta$ quedan determinados por
$f$ de manera única, y cada pareja $(\alpha,\beta)$ corresponde a una única función de $V$.
Del mismo modo, las coordenadas complejas de $g$ y $h$ se denotan
$\alpha'=a_1'+ib_1'$, $\beta'=a_2'+ib_2'$ y $\alpha''=a_1''+ib_1''$, $\beta''=a_2''+ib_2''$.

La suma también se describe en estas coordenadas. Por $f+g=(a_1+a_1')\cos t+(b_1+b_1')\sin t+(a_2+a_2')\cos 2t+(b_2+b_2')\sin 2t$ \eqref{eq:suma}, los coeficientes de
$\cos t$ y $\sin t$ en $f+g$ son $a_1+a_1'$ y $b_1+b_1'$, así que su primera coordenada
compleja es $(a_1+a_1')+i(b_1+b_1')=\alpha+\alpha'$; de la misma forma, la segunda es
$\beta+\beta'$. Es decir,
\begin{equation}\label{eq:suma-coord}
  f+g=(\alpha+\alpha')\cdot\cos t+(\beta+\beta')\cdot\cos 2t ,
\end{equation}
y sumar en $V$ equivale a sumar las coordenadas complejas una a una.

\paragraph{Definición de la multiplicación.}
Para $z\in\mathbb C$ y $f=\alpha\cdot\cos t+\beta\cdot\cos 2t$ se define
\begin{equation}\label{eq:mult}
  \destacado{z\cdot f=(z\alpha)\cdot\cos t+(z\beta)\cdot\cos 2t ,}
\end{equation}
es decir, multiplicar $f$ por $z$ consiste en multiplicar por $z$ cada una de sus coordenadas
complejas. La definición no presenta ambigüedad, porque $\alpha$ y $\beta$ son únicos, y es
compatible con el enunciado: para $f=\cos t$ se tiene $\alpha=1$, $\beta=0$, y
$z\cdot f=(z\alpha)\cdot\cos t+(z\beta)\cdot\cos 2t$ \eqref{eq:mult} se reduce a $z\cdot\cos t$ tal como lo da la regla; lo mismo ocurre con
$\cos 2t$. Además, el resultado pertenece a $V$: si $z\alpha=p+iq$ con $p,q\in\mathbb R$, por
$(a+ib)\cdot\cos nt=a\cos nt+b\sin nt$ \eqref{eq:regla-coef} se tiene $(z\alpha)\cdot\cos t=p\cos t+q\sin t$, y análogamente para
$z\beta$, de modo que $z\cdot f$ es una combinación con coeficientes reales de las cuatro
funciones generadoras.

\paragraph{Multiplicación por $i$.}
Este caso se usa de nuevo en la construcción del producto hermítico. Como
$i\alpha=i(a_1+ib_1)=-b_1+ia_1$ e $i\beta=-b_2+ia_2$, la definición $z\cdot f=(z\alpha)\cdot\cos t+(z\beta)\cdot\cos 2t$ \eqref{eq:mult} junto con
$(a+ib)\cdot\cos nt=a\cos nt+b\sin nt$ \eqref{eq:regla-coef} da
\begin{equation}\label{eq:i-f}
  \destacado{i\cdot f=-b_1\cos t+a_1\sin t-b_2\cos 2t+a_2\sin 2t .}
\end{equation}
En particular, $\sin nt$ tiene coordenada compleja $i$ (porque $\sin nt=i\cdot\cos nt$ por
$i\cdot\cos nt=\sin nt$ \eqref{eq:i-cos}), así que $i\cdot\sin nt=(i\cdot i)\cdot\cos nt=-\cos nt$. Multiplicar por
$i$ desplaza cada onda un cuarto de período: $\cos nt$ pasa a $\sin nt$ y $\sin nt$ pasa a
$-\cos nt$.

\paragraph{(M1) $z\cdot(f+g)=z\cdot f+z\cdot g$.}
Por $f+g=(\alpha+\alpha')\cdot\cos t+(\beta+\beta')\cdot\cos 2t$ \eqref{eq:suma-coord}, $f+g$ tiene coordenadas $\alpha+\alpha'$ y $\beta+\beta'$, y por
$z\cdot f=(z\alpha)\cdot\cos t+(z\beta)\cdot\cos 2t$ \eqref{eq:mult},
\[
  z\cdot(f+g)=\bigl(z(\alpha+\alpha')\bigr)\cdot\cos t
  +\bigl(z(\beta+\beta')\bigr)\cdot\cos 2t .
\]
Por otra parte, $z\cdot f$ tiene coordenadas $z\alpha$, $z\beta$ y $z\cdot g$ tiene
coordenadas $z\alpha'$, $z\beta'$, de modo que, por $f+g=(\alpha+\alpha')\cdot\cos t+(\beta+\beta')\cdot\cos 2t$ \eqref{eq:suma-coord},
\[
  z\cdot f+z\cdot g=(z\alpha+z\alpha')\cdot\cos t+(z\beta+z\beta')\cdot\cos 2t .
\]
La multiplicación de números complejos se distribuye sobre la suma, así que
$z(\alpha+\alpha')=z\alpha+z\alpha'$ y $z(\beta+\beta')=z\beta+z\beta'$. Ambos lados tienen las
mismas coordenadas y, por tanto, $z\cdot(f+g)=z\cdot f+z\cdot g$.

\paragraph{(M2) $(z+w)\cdot f=z\cdot f+w\cdot f$.}
Por $z\cdot f=(z\alpha)\cdot\cos t+(z\beta)\cdot\cos 2t$ \eqref{eq:mult}, las coordenadas de $(z+w)\cdot f$ son $(z+w)\alpha$ y $(z+w)\beta$. Por
$z\cdot f=(z\alpha)\cdot\cos t+(z\beta)\cdot\cos 2t$ \eqref{eq:mult} y $f+g=(\alpha+\alpha')\cdot\cos t+(\beta+\beta')\cdot\cos 2t$ \eqref{eq:suma-coord}, las de $z\cdot f+w\cdot f$ son $z\alpha+w\alpha$ y
$z\beta+w\beta$. Por la propiedad distributiva en $\mathbb C$, $(z+w)\alpha=z\alpha+w\alpha$ y
$(z+w)\beta=z\beta+w\beta$, de modo que las coordenadas coinciden y
$(z+w)\cdot f=z\cdot f+w\cdot f$.

\paragraph{(M3) $(zw)\cdot f=z\cdot(w\cdot f)$.}
Las coordenadas de $(zw)\cdot f$ son $(zw)\alpha$ y $(zw)\beta$. Por otro lado, $w\cdot f$ tiene
coordenadas $w\alpha$ y $w\beta$, y al multiplicar por $z$ se obtienen $z(w\alpha)$ y
$z(w\beta)$. Como la multiplicación de números complejos es asociativa,
$(zw)\alpha=z(w\alpha)$ y $(zw)\beta=z(w\beta)$, de donde $(zw)\cdot f=z\cdot(w\cdot f)$.

\paragraph{(M4) $1\cdot f=f$.}
Con $z=1$, la definición $z\cdot f=(z\alpha)\cdot\cos t+(z\beta)\cdot\cos 2t$ \eqref{eq:mult} da
$1\cdot f=(1\alpha)\cdot\cos t+(1\beta)\cdot\cos 2t=\alpha\cdot\cos t+\beta\cdot\cos 2t=f$.

\medskip
Así, la multiplicación por escalares complejos produce elementos de $V$ y se cumplen (M1)--(M4). Junto
con lo demostrado para la suma, esto prueba que $V$ es un espacio vectorial complejo.

%======================================================================
\subsection*{Dimensión compleja de $V$}
%======================================================================

\begin{definicionT}[Base y dimensión]
Un conjunto $\{v_1,\dots,v_k\}$ de vectores de un espacio vectorial complejo $V$ es una
\emph{base} de $V$ si cumple dos condiciones. La primera es que \emph{genera} $V$: todo
vector de $V$ se escribe como $c_1\cdot v_1+\dots+c_k\cdot v_k$ con $c_1,\dots,c_k\in\mathbb C$.
La segunda es que es \emph{linealmente independiente sobre $\mathbb C$}: si
$c_1\cdot v_1+\dots+c_k\cdot v_k=0$ con $c_j\in\mathbb C$, entonces $c_1=\dots=c_k=0$. Todas
las bases de un mismo espacio tienen el mismo número de elementos, y ese número se llama
\emph{dimensión} del espacio; para un espacio complejo se denota $\dim_{\mathbb C}V$.
\end{definicionT}

El conjunto $\{\cos t,\cos 2t\}$ genera $V$, pues por $f=\alpha\cdot\cos t+\beta\cdot\cos 2t$ \eqref{eq:alfa-beta} todo elemento de $V$
se escribe como $\alpha\cdot\cos t+\beta\cdot\cos 2t$ con $\alpha,\beta\in\mathbb C$. Para la
independencia, supóngase que $\alpha\cdot\cos t+\beta\cdot\cos 2t=0$ con
$\alpha=a_1+ib_1$ y $\beta=a_2+ib_2$. Por $(a+ib)\cdot\cos nt=a\cos nt+b\sin nt$ \eqref{eq:regla-coef}, el lado izquierdo es la función
$a_1\cos t+b_1\sin t+a_2\cos 2t+b_2\sin 2t$, que es la función cero; por el
lema~\ref{lema:independencia}, $a_1=b_1=a_2=b_2=0$, es decir, $\alpha=\beta=0$. En
consecuencia, $\{\cos t,\cos 2t\}$ es una base de $V$ y
\[
  \destacado{\dim_{\mathbb C}V=2 .}
\]

\begin{definicionT}[Notación de Dirac]
En computación cuántica, los vectores del espacio de estados se escriben con el símbolo
$\ket{\psi}$, llamado \emph{ket}, donde $\psi$ es solo un nombre. Siguiendo el enunciado, se
denota
\[
  \ket0=\cos t,\qquad \ket1=\cos 2t .
\]
\end{definicionT}

Con esta notación, $f=\alpha\cdot\cos t+\beta\cdot\cos 2t$ \eqref{eq:alfa-beta} se escribe $f=\alpha\ket0+\beta\ket1$, y el resultado
anterior dice que $V=\spanC\{\ket0,\ket1\}$, donde $\spanC$ indica combinaciones lineales con
coeficientes complejos.


%======================================================================
\subsection*{El producto interno de $L^2$ en $V$}
%======================================================================

\begin{definicionT}[El espacio $L^2$ y su producto interno]
El espacio $L^2([0,2\pi])$ está formado por las funciones definidas en $[0,2\pi]$ cuyo
cuadrado tiene integral finita en ese intervalo. Para dos funciones reales $f,g$ de este
espacio, su \emph{producto interno de $L^2$} es
\[
  \destacado{\Lp fg=\int_0^{2\pi}f(t)\,g(t)\,dt .}
\]
Para funciones con valores complejos aparece además un conjugado en uno de los factores, pero
todos los elementos de $V$ son funciones reales, de modo que ese conjugado no interviene aquí.
El subíndice $L^2$ indica el espacio del que proviene el producto, y los paréntesis lo
distinguen del producto hermítico $\bk fg$ que se construye más adelante.
\end{definicionT}

Los elementos de $V$ pertenecen a $L^2([0,2\pi])$ porque son funciones continuas en un
intervalo cerrado y acotado. El producto $\Lp fg$ es un número real y es simétrico,
$\Lp fg=\Lp gf$, porque $f(t)g(t)=g(t)f(t)$. Además, como la integral de una suma es la suma de
las integrales y las constantes salen de la integral, es lineal en cada argumento: para
$c,d\in\mathbb R$,
\begin{equation}\label{eq:L2-lineal}
  \Lp{cf+dg}{h}=c\,\Lp fh+d\,\Lp gh,
  \qquad
  \Lp{f}{cg+dh}=c\,\Lp fg+d\,\Lp fh .
\end{equation}
Las constantes que salen de este producto son números \emph{reales}, y esa restricción es la
que obliga más adelante a modificarlo.

\paragraph{Reducción a productos de senos y cosenos.}
Para calcular $\Lp fg$ con $f$ y $g$ arbitrarios en $V$ se vuelve a su escritura
$f(t)=a_1\cos t+b_1\sin t+a_2\cos 2t+b_2\sin 2t$ \eqref{eq:forma-general}. Para abreviar, las cuatro funciones generadoras se nombran
\[
  \varphi_1=\Cu,\qquad \varphi_2=\Su,\qquad \varphi_3=\Cd,\qquad \varphi_4=\Sd,
\]
y los coeficientes de $f$ y $g$ se agrupan como
$(c_1,c_2,c_3,c_4)=(a_1,b_1,a_2,b_2)$ y $(c_1',c_2',c_3',c_4')=(a_1',b_1',a_2',b_2')$, de modo
que $f=\sum_{j=1}^4 c_j\varphi_j$ y $g=\sum_{k=1}^4 c_k'\varphi_k$. Por la linealidad
$(cf+dg,h)_{L^2}=c\,(f,h)_{L^2}+d\,(g,h)_{L^2},\quad\penalty0 (f,cg+dh)_{L^2}=c\,(f,g)_{L^2}+d\,(f,h)_{L^2}$ \eqref{eq:L2-lineal}, los coeficientes salen del producto:
\begin{equation}\label{eq:L2-expansion}
  \Lp fg=\int_0^{2\pi}\Bigl(\sum_{j=1}^4 c_j\varphi_j(t)\Bigr)\Bigl(\sum_{k=1}^4 c_k'\varphi_k(t)\Bigr)dt
  =\sum_{j=1}^4\sum_{k=1}^4 c_j\,c_k'\,\Lp{\varphi_j}{\varphi_k} .
\end{equation}
Calcular el producto de dos elementos cualesquiera de $V$ se reduce entonces a conocer los
dieciséis productos $\Lp{\varphi_j}{\varphi_k}$ entre las cuatro funciones con las que se
construyen $f$ y $g$. Cada uno es la integral de un producto de dos funciones
trigonométricas, de la forma $\cos(nt)\cos(mt)$, $\sin(nt)\sin(mt)$ o $\cos(nt)\sin(mt)$ con
$n,m\in\{1,2\}$; esa es la razón por la que a continuación se trabaja con productos de senos y
cosenos.

\paragraph{Identidades de producto a suma.}
Un producto de dos funciones trigonométricas no se integra de forma directa, pero puede
reescribirse como una suma de un solo coseno o un solo seno, cuya integral es inmediata. Para
ello se parte de las fórmulas de suma y resta de ángulos,
\begin{align*}
\cos(A-B) &= \cos A\cos B + \sin A\sin B, & \cos(A+B) &= \cos A\cos B - \sin A\sin B,\\
\sin(A+B) &= \sin A\cos B + \cos A\sin B, & \sin(A-B) &= \sin A\cos B - \cos A\sin B.
\end{align*}
Al sumar las dos primeras se cancela $\sin A\sin B$; al restarlas se cancela
$\cos A\cos B$; y al restar la cuarta de la tercera se cancela $\sin A\cos B$. Despejando, y
tomando $A=nt$, $B=mt$, se obtiene
\begin{equation}\label{eq:producto-suma}
\begin{aligned}
\cos(nt)\cos(mt) &= \tfrac12\bigl[\cos((n-m)t) + \cos((n+m)t)\bigr],\\
\sin(nt)\sin(mt) &= \tfrac12\bigl[\cos((n-m)t) - \cos((n+m)t)\bigr],\\
\cos(nt)\sin(mt) &= \tfrac12\bigl[\sin((n+m)t) - \sin((n-m)t)\bigr].
\end{aligned}
\end{equation}
Después de aplicar estas identidades, solo quedan por calcular integrales de la forma
$\int_0^{2\pi}\cos(kt)\,dt$ y $\int_0^{2\pi}\sin(kt)\,dt$ con $k$ entero.

\paragraph{Integral de un solo coseno o un solo seno.}
La figura~\ref{fig:integrales} muestra por qué estas integrales valen cero cuando $k\ne 0$. En
$[0,2\pi]$, la función $\cos(kt)$ o $\sin(kt)$ completa exactamente $k$ oscilaciones, y en cada
oscilación el tramo de curva que queda por encima del eje encierra la misma área que el que
queda por debajo, pues uno es el reflejo del otro. Los ceros de la función separan esos tramos:
$\cos t$ se anula en $\pi/2$ y $3\pi/2$; $\cos 2t$ en $\pi/4$, $3\pi/4$, $5\pi/4$ y $7\pi/4$;
$\sin t$ en $0$, $\pi$ y $2\pi$; y $\sin 2t$ en $0$, $\pi/2$, $\pi$, $3\pi/2$ y $2\pi$. Como la
integral cuenta con signo positivo el área por encima del eje y con signo negativo el área
por debajo, las dos contribuciones se cancelan.

\begin{figure}[H]
\centering
\panelarea{cos(x)}{colCu}{$\cos t$}{(1.5708,0) (4.7124,0)}{
  \node[font=\small] at (axis cs:0.6,0.35) {$+$};
  \node[font=\small] at (axis cs:3.1416,-0.4) {$-$};
  \node[font=\small] at (axis cs:5.7,0.35) {$+$};}
\hfill
\panelarea{cos(2*x)}{colCd}{$\cos 2t$}{(0.7854,0) (2.3562,0) (3.9270,0) (5.4978,0)}{
  \node[font=\small] at (axis cs:1.5708,-0.4) {$-$};
  \node[font=\small] at (axis cs:3.1416,0.4) {$+$};
  \node[font=\small] at (axis cs:4.7124,-0.4) {$-$};}

\vspace{2mm}
\panelarea{sin(x)}{colSu}{$\sin t$}{(0,0) (3.1416,0) (6.2832,0)}{
  \node[font=\small] at (axis cs:1.5708,0.4) {$+$};
  \node[font=\small] at (axis cs:4.7124,-0.4) {$-$};}
\hfill
\panelarea{sin(2*x)}{colSd}{$\sin 2t$}{(0,0) (1.5708,0) (3.1416,0) (4.7124,0) (6.2832,0)}{
  \node[font=\small] at (axis cs:0.7854,0.4) {$+$};
  \node[font=\small] at (axis cs:2.3562,-0.4) {$-$};
  \node[font=\small] at (axis cs:3.9270,0.4) {$+$};
  \node[font=\small] at (axis cs:5.4978,-0.4) {$-$};}
\caption{Cosenos y senos de frecuencia $1$ y $2$ en $[0,2\pi]$. Los puntos rojos son los ceros
de cada función. El área en azul (por encima del eje) es igual al área en rojo (por debajo),
así que la integral de cada función sobre $[0,2\pi]$ es $0$.}
\label{fig:integrales}
\end{figure}

El cálculo confirma lo que muestra la figura. Si $k$ es un entero distinto de cero, una
antiderivada de $\cos(kt)$ es $\sin(kt)/k$ y una de $\sin(kt)$ es $-\cos(kt)/k$, de modo que
\begin{align*}
  \int_0^{2\pi}\cos(kt)\,dt&=\left[\frac{\sin(kt)}{k}\right]_0^{2\pi}
  =\frac{\sin(2\pi k)-\sin 0}{k}=0,\\[1mm]
  \int_0^{2\pi}\sin(kt)\,dt&=\left[-\frac{\cos(kt)}{k}\right]_0^{2\pi}
  =\frac{\cos 0-\cos(2\pi k)}{k}=0,
\end{align*}
ya que, por ser $2\pi k$ un múltiplo entero de $2\pi$, $\sin(2\pi k)=0$ y
$\cos(2\pi k)=1=\cos 0$. Si $k=0$, el coseno es la función constante $\cos(0\cdot t)=1$, cuya
gráfica encierra con el eje un rectángulo de base $2\pi$ y altura $1$, mientras que el seno es
la función constante $0$. En resumen,
\begin{equation}\label{eq:integrales}
  \destacado{\int_0^{2\pi}\cos(kt)\,dt=
  \begin{cases} 2\pi, & k=0,\\ 0, & k\ne 0,\end{cases}
  \qquad
  \int_0^{2\pi}\sin(kt)\,dt=0\ \text{ para todo entero } k.}
\end{equation}

\paragraph{Productos entre las funciones generadoras.}
Las identidades \eqref{eq:producto-suma} son las siguientes:
\begin{align*}
\cos nt\cos mt&=\tfrac12\bigl[\cos((n-m)t)+\cos((n+m)t)\bigr],\\
\sin nt\sin mt&=\tfrac12\bigl[\cos((n-m)t)-\cos((n+m)t)\bigr],\\
\cos nt\sin mt&=\tfrac12\bigl[\sin((n+m)t)-\sin((n-m)t)\bigr]
\end{align*}
Al integrarlas y aplicar $\int_0^{2\pi}\cos(kt)\,dt=\begin{cases}2\pi,&k=0\\0,&k\neq0\end{cases}\quad\text{y}\quad \int_0^{2\pi}\sin(kt)\,dt=0$ \eqref{eq:integrales}, el valor
de cada producto depende de si las frecuencias $n-m$ y $n+m$ son o no cero. Como
$n,m\in\{1,2\}$, la frecuencia $n+m$ vale $2$, $3$ o $4$ y nunca es cero, mientras que $n-m$ es
cero exactamente cuando $n=m$. Para dos cosenos,
\[
  \Lp{\cos nt}{\cos mt}
  =\frac12\int_0^{2\pi}\cos((n-m)t)\,dt+\frac12\underbrace{\int_0^{2\pi}\cos((n+m)t)\,dt}_{=\,0},
\]
y el primer término vale $\frac12\cdot 2\pi=\pi$ si $n=m$ y $0$ si $n\ne m$. Para dos senos
ocurre lo mismo, pues la identidad tiene los mismos dos cosenos con un signo menos en el
segundo, que de todos modos se anula. Para un coseno y un seno, la identidad solo contiene
senos, y por $\int_0^{2\pi}\cos(kt)\,dt=\begin{cases}2\pi,&k=0\\0,&k\neq0\end{cases}\quad\text{y}\quad \int_0^{2\pi}\sin(kt)\,dt=0$ \eqref{eq:integrales} el resultado es $0$ para cualquier $n$ y $m$.

La figura~\ref{fig:productos} ilustra estos resultados. La función
$\cos^2 t=\frac12+\frac12\cos 2t$ nunca es negativa y oscila alrededor de su valor medio
$\frac12$, de modo que el área bajo su gráfica es la de un rectángulo de base $2\pi$ y altura
$\frac12$, es decir, $\pi$. En cambio, el producto de dos funciones distintas toma valores
positivos y negativos que se compensan exactamente.

\begin{figure}[H]
\centering
\panelarea{cos(x)^2}{colCu}{$\cos t\cdot\cos t$ \quad (integral $=\pi$)}{(1.5708,0) (4.7124,0)}{
  \draw[dashed, thick, resaltado] (axis cs:0,0.5) -- (axis cs:6.2832,0.5);
  \node[resaltado, font=\footnotesize, anchor=west] at (axis cs:6.35,0.5) {$\tfrac12$};}
\hfill
\panelarea{cos(x)*cos(2*x)}{black!75}{$\cos t\cdot\cos 2t$ \quad (integral $=0$)}%
  {(0.7854,0) (1.5708,0) (2.3562,0) (3.9270,0) (4.7124,0) (5.4978,0)}{}

\vspace{2mm}
\panelarea{cos(x)*sin(x)}{black!75}{$\cos t\cdot\sin t$ \quad (integral $=0$)}%
  {(0,0) (1.5708,0) (3.1416,0) (4.7124,0) (6.2832,0)}{}
\hfill
\panelarea{sin(x)*sin(2*x)}{black!75}{$\sin t\cdot\sin 2t$ \quad (integral $=0$)}%
  {(0,0) (1.5708,0) (3.1416,0) (4.7124,0) (6.2832,0)}{}
\caption{Productos entre funciones generadoras. Arriba a la izquierda, una función por sí
misma: toda el área queda por encima del eje y la integral es $2\pi\cdot\frac12=\pi$. En los
otros tres paneles, dos funciones distintas: las áreas azul y roja son iguales y la integral
es $0$.}
\label{fig:productos}
\end{figure}

\begin{definicionT}[Delta de Kronecker]
Para enteros $j,k$ se define
\[
  \delta_{jk}=\begin{cases} 1, & j=k,\\ 0, & j\ne k.\end{cases}
\]
Esta notación permite escribir en una sola fórmula los casos ``índices iguales'' e ``índices
distintos''.
\end{definicionT}

Con ella, los resultados anteriores se resumen en
\begin{equation}\label{eq:ortogonalidad}
  \Lp{\cos nt}{\cos mt}=\pi\,\delta_{nm},\qquad
  \Lp{\sin nt}{\sin mt}=\pi\,\delta_{nm},\qquad
  \Lp{\cos nt}{\sin mt}=0,
\end{equation}
o, en términos de las funciones $\varphi_j$, $\Lp{\varphi_j}{\varphi_k}=\pi\,\delta_{jk}$. La
tabla siguiente reúne los dieciséis productos; en la fila $j$ y la columna $k$ aparece
$\Lp{\varphi_j}{\varphi_k}$.
\[
\renewcommand{\arraystretch}{1.3}
\begin{array}{c|cccc}
\Lp{\cdot}{\cdot} & \Cu & \Su & \Cd & \Sd\\ \hline
\Cu & \color{resaltado}\boldsymbol\pi & \color{ceroCol}0 & \color{ceroCol}0 & \color{ceroCol}0\\
\Su & \color{ceroCol}0 & \color{resaltado}\boldsymbol\pi & \color{ceroCol}0 & \color{ceroCol}0\\
\Cd & \color{ceroCol}0 & \color{ceroCol}0 & \color{resaltado}\boldsymbol\pi & \color{ceroCol}0\\
\Sd & \color{ceroCol}0 & \color{ceroCol}0 & \color{ceroCol}0 & \color{resaltado}\boldsymbol\pi
\end{array}
\]
Los valores $\pi$ ocupan la diagonal, donde cada función se multiplica por sí misma; fuera de
ella todos los productos son $0$, es decir, las cuatro funciones son \emph{ortogonales} entre sí
respecto a $\Lp{\cdot}{\cdot}$.

\paragraph{Fórmula del producto en términos de los coeficientes.}
En la suma doble \eqref{eq:L2-expansion}, que es
\begin{equation*}
(f,g)_{L^2}=\sum_{j=1}^4\sum_{k=1}^4 c_j\,c_k'\,(\varphi_j,\varphi_k)_{L^2},
\end{equation*}
todos los términos con $j\ne k$ se anulan y solo
sobreviven los cuatro términos con $j=k$, cada uno multiplicado por $\pi$. Por tanto,
\begin{equation}\label{eq:L2-formula}
  \destacado{\Lp fg=\pi\bigl(a_1a_1'+b_1b_1'+a_2a_2'+b_2b_2'\bigr).}
\end{equation}
El producto de $L^2$ en $V$ es, salvo el factor $\pi$, la suma de los productos de
coeficientes correspondientes, igual que el producto punto de dos vectores de
$\mathbb R^4$.

%======================================================================
\subsection*{Producto hermítico}
%======================================================================

\begin{definicionT}[Producto hermítico]
Un \emph{producto hermítico} en un espacio vectorial complejo $V$ es una regla que asigna a
cada par de vectores $f,g$ un número complejo $\bk fg$ y que cumple, para todo $f,g,h\in V$ y
todo $z,w\in\mathbb C$, las propiedades siguientes.
\begin{enumerate}
  \item[(H1)] \emph{Hermiticidad:} $\bk gf=\overline{\bk fg}$.
  \item[(H2)] \emph{Linealidad en el segundo argumento:}
  $\bk{f}{z\cdot g+w\cdot h}=z\,\bk fg+w\,\bk fh$.
  \item[(H3)] \emph{Antilinealidad en el primer argumento:}
  $\bk{z\cdot f+w\cdot g}{h}=\overline z\,\bk fh+\overline w\,\bk gh$.
  \item[(H4)] \emph{Positividad definida:} $\bk ff\ge 0$, y $\bk ff=0$ solo si $f=0$.
\end{enumerate}
\end{definicionT}

\paragraph{Por qué $\Lp{\cdot}{\cdot}$ no es un producto hermítico.}
El producto de $L^2$ siempre da números reales y solo permite sacar escalares reales, como
indica $(cf+dg,h)_{L^2}=c\,(f,h)_{L^2}+d\,(g,h)_{L^2},\quad\penalty0 (f,cg+dh)_{L^2}=c\,(f,g)_{L^2}+d\,(f,h)_{L^2}$ \eqref{eq:L2-lineal}. Con escalares complejos falla la propiedad (H2). Por ejemplo, por
$i\cdot\cos nt=\sin nt$ \eqref{eq:i-cos}, $i\cdot\cos t=\sin t$, así que
\[
  \Lp{\cos t}{i\cdot\cos t}=\Lp{\cos t}{\sin t}=0 .
\]
Si $\Lp{\cdot}{\cdot}$ cumpliera (H2), debería tenerse
\[
  \Lp{\cos t}{i\cdot\cos t}=i\,\Lp{\cos t}{\cos t}=i\pi\ne 0 ,
\]
lo cual contradice el cálculo anterior. El problema de fondo es que $\Lp{\cdot}{\cdot}$ trata a
$\cos t$ y $\sin t$ como funciones perpendiculares, mientras que en $V$, visto como espacio
complejo, $\sin t$ es simplemente $i$ veces $\cos t$. Un producto que solo ve coeficientes
reales no puede detectar esa relación, y por eso pierde la información de la parte imaginaria.

\paragraph{El producto buscado en coordenadas complejas.}
Como $\ket0$ y $\ket1$ tienen coordenadas complejas $(1,0)$ y $(0,1)$, un producto hermítico en
el que $\{\ket0,\ket1\}$ sea ortonormal debe comportarse, en coordenadas, como el producto
hermítico usual de $\mathbb C^2$. Para $f=\alpha\ket0+\beta\ket1$ y
$g=\alpha'\ket0+\beta'\ket1$, ese producto es
\begin{equation}\label{eq:producto-C2}
  \bk fg=\overline\alpha\,\alpha'+\overline\beta\,\beta' .
\end{equation}
El enunciado pide construirlo a partir de $\Lp{\cdot}{\cdot}$, de modo que el objetivo es
escribir $\bk fg=\overline\alpha\,\alpha'+\overline\beta\,\beta'$ \eqref{eq:producto-C2} en términos de ese producto. Al sustituir
$\alpha=a_1+ib_1$, $\beta=a_2+ib_2$, $\alpha'=a_1'+ib_1'$ y $\beta'=a_2'+ib_2'$ y desarrollar,
\[
\begin{aligned}
\overline\alpha\,\alpha'+\overline\beta\,\beta'
&=(a_1-ib_1)(a_1'+ib_1')+(a_2-ib_2)(a_2'+ib_2')\\
&=\bigl(a_1a_1'+b_1b_1'+a_2a_2'+b_2b_2'\bigr)
 +i\bigl(a_1b_1'-b_1a_1'+a_2b_2'-b_2a_2'\bigr).
\end{aligned}
\]
Por $(f,g)_{L^2}=\pi\bigl(a_1a_1'+b_1b_1'+a_2a_2'+b_2b_2'\bigr)$ \eqref{eq:L2-formula}, la parte real es exactamente $\frac1\pi\Lp fg$. Falta expresar la
parte imaginaria.

\paragraph{La parte imaginaria a partir de la multiplicación por $i$.}
Por $i\cdot f=-b_1\cos t+a_1\sin t-b_2\cos 2t+a_2\sin 2t$ \eqref{eq:i-f}, los coeficientes de $i\cdot f$ son $(-b_1,\ a_1,\ -b_2,\ a_2)$. Al aplicar
$(f,g)_{L^2}=\pi\bigl(a_1a_1'+b_1b_1'+a_2a_2'+b_2b_2'\bigr)$ \eqref{eq:L2-formula} a $i\cdot f$ y $g$ se obtiene
\[
  \Lp{i\cdot f}{g}=\pi\bigl(-b_1a_1'+a_1b_1'-b_2a_2'+a_2b_2'\bigr),
\]
que es $\pi$ veces la parte imaginaria hallada arriba. La multiplicación por $i$ intercambia
los papeles de coseno y seno, y eso es justamente lo que permite que $\Lp{\cdot}{\cdot}$ capture
los productos cruzados que forman la parte imaginaria.

\paragraph{Definición del producto hermítico.}
Por lo anterior, se define, para $f,g\in V$,
\begin{equation}\label{eq:hermitico}
  \destacado{\bk fg=\frac1\pi\Bigl[\Lp fg+i\,\Lp{i\cdot f}{g}\Bigr].}
\end{equation}
El primer sumando aporta la parte real y el segundo, multiplicado por $i$, la parte
imaginaria; el factor $1/\pi$ compensa el $\pi$ que aparece en $(f,g)_{L^2}=\pi\bigl(a_1a_1'+b_1b_1'+a_2a_2'+b_2b_2'\bigr)$ \eqref{eq:L2-formula} y es lo
que hace que $\ket0$ y $\ket1$ tengan norma $1$. Por los dos cálculos anteriores,
\begin{equation}\label{eq:hermitico-coord}
  \destacado{\bk fg=\overline\alpha\,\alpha'+\overline\beta\,\beta' ,}
\end{equation}
es decir, en coordenadas complejas el producto construido coincide con el producto usual de
$\mathbb C^2$. Esta expresión es la que se usa para verificar las propiedades.

%======================================================================
\subsection*{Verificación de las propiedades del producto hermítico}
%======================================================================

Sean $f=\alpha\ket0+\beta\ket1$, $g=\alpha'\ket0+\beta'\ket1$ y $h=\alpha''\ket0+\beta''\ket1$
elementos de $V$, y sean $z,w\in\mathbb C$.

\paragraph{(H1) Hermiticidad.}
Por $\bk fg=\overline\alpha\,\alpha'+\overline\beta\,\beta'$ \eqref{eq:hermitico-coord} y las propiedades del conjugado $\overline{z+w}=\overline z+\overline w,\quad\penalty0 \overline{zw}=\overline z\,\overline w,\quad\penalty0 \overline{\overline z}=z$ \eqref{eq:conjugado},
\[
  \overline{\bk fg}
  =\overline{\overline\alpha\,\alpha'+\overline\beta\,\beta'}
  =\overline{\overline\alpha}\;\overline{\alpha'}+\overline{\overline\beta}\;\overline{\beta'}
  =\alpha\,\overline{\alpha'}+\beta\,\overline{\beta'}
  =\overline{\alpha'}\,\alpha+\overline{\beta'}\,\beta
  =\bk gf ,
\]
donde en el penúltimo paso se usó que el producto de números complejos no depende del orden,
y en el último, que $\bk gf$ se obtiene de $\bk fg=\overline\alpha\,\alpha'+\overline\beta\,\beta'$ \eqref{eq:hermitico-coord} conjugando las
coordenadas de $g$, que ahora está en el primer argumento.

\paragraph{(H2) Linealidad en el segundo argumento.}
Para usar $\bk fg=\overline\alpha\,\alpha'+\overline\beta\,\beta'$ \eqref{eq:hermitico-coord} con $z\cdot g+w\cdot h$ en el segundo argumento se
necesitan sus coordenadas complejas. Por la definición $z\cdot f=(z\alpha)\cdot\cos t+(z\beta)\cdot\cos 2t$ \eqref{eq:mult}, $z\cdot g$ tiene
coordenadas $z\alpha'$ y $z\beta'$, y $w\cdot h$ tiene coordenadas $w\alpha''$ y $w\beta''$;
por $f+g=(\alpha+\alpha')\cdot\cos t+(\beta+\beta')\cdot\cos 2t$ \eqref{eq:suma-coord}, al sumarlas se suman las coordenadas, de modo que
\[
  z\cdot g+w\cdot h=(z\alpha'+w\alpha'')\ket0+(z\beta'+w\beta'')\ket1 .
\]
Entonces
\[
\begin{aligned}
  \bk{f}{z\cdot g+w\cdot h}
  &=\overline\alpha\,(z\alpha'+w\alpha'')+\overline\beta\,(z\beta'+w\beta'')\\
  &=z\bigl(\overline\alpha\,\alpha'+\overline\beta\,\beta'\bigr)
   +w\bigl(\overline\alpha\,\alpha''+\overline\beta\,\beta''\bigr)
   =z\,\bk fg+w\,\bk fh ,
\end{aligned}
\]
donde se distribuyeron los productos y se agruparon los términos que contienen $z$ y los que
contienen $w$. Los dos paréntesis son, por $\bk fg=\overline\alpha\,\alpha'+\overline\beta\,\beta'$ \eqref{eq:hermitico-coord}, $\bk fg$ y $\bk fh$.

\paragraph{(H3) Antilinealidad en el primer argumento.}
Esta propiedad se deduce de las dos anteriores. Por (H1) se pasa la combinación al segundo
argumento, por (H2) se sacan los escalares y por $\overline{z+w}=\overline z+\overline w,\quad\penalty0 \overline{zw}=\overline z\,\overline w,\quad\penalty0 \overline{\overline z}=z$ \eqref{eq:conjugado} se distribuye el
conjugado:
\[
\begin{aligned}
  \bk{z\cdot f+w\cdot g}{h}
  &=\overline{\bk{h}{z\cdot f+w\cdot g}}
   =\overline{z\,\bk hf+w\,\bk hg}\\
  &=\overline z\;\overline{\bk hf}+\overline w\;\overline{\bk hg}
   =\overline z\,\bk fh+\overline w\,\bk gh ,
\end{aligned}
\]
donde en el último paso se usó de nuevo (H1).

\paragraph{(H4) Positividad definida.}
Por $\bk fg=\overline\alpha\,\alpha'+\overline\beta\,\beta'$ \eqref{eq:hermitico-coord} y $z\,\overline z=|z|^2$ \eqref{eq:modulo},
\[
  \bk ff=\overline\alpha\,\alpha+\overline\beta\,\beta=|\alpha|^2+|\beta|^2
  =a_1^2+b_1^2+a_2^2+b_2^2\ \ge 0 ,
\]
porque es una suma de cuadrados de números reales. Además, una suma de cuadrados de números
reales es cero solo si cada sumando es cero, así que $\bk ff=0$ implica
$a_1=b_1=a_2=b_2=0$, es decir, que $f$ es la función cero. Recíprocamente, si $f=0$, el
lema~\ref{lema:independencia} garantiza que sus cuatro coeficientes son cero y, por tanto,
$\bk ff=0$.

\medskip
Con esto, $\bk{\cdot}{\cdot}$ es un producto hermítico en $V$.

%======================================================================
\subsection*{Ortonormalidad de $\{\ket0,\ket1\}$}
%======================================================================

\begin{definicionT}[Norma, ortogonalidad y base ortonormal]
En un espacio con producto hermítico, la \emph{norma} de un vector es
$\lVert f\rVert=\sqrt{\bk ff}$, que mide su longitud. Dos vectores $f,g$ son
\emph{ortogonales} si $\bk fg=0$. Una base es \emph{ortonormal} si cada uno de sus elementos
tiene norma $1$ y sus elementos son ortogonales entre sí.
\end{definicionT}

Para $\{\ket0,\ket1\}$ esto significa comprobar que $\bk00=\bk11=1$ y $\bk01=0$; la igualdad
$\bk10=0$ se sigue entonces de (H1), pues $\bk10=\overline{\bk01}=0$. Las coordenadas complejas
de $\ket0=\cos t$ son $(1,0)$ y las de $\ket1=\cos 2t$ son $(0,1)$, de modo que, por
$\bk fg=\overline\alpha\,\alpha'+\overline\beta\,\beta'$ \eqref{eq:hermitico-coord},
\[
  \bk00=\overline1\cdot1+\overline0\cdot0=1,\qquad
  \bk11=\overline0\cdot0+\overline1\cdot1=1,\qquad
  \bk01=\overline1\cdot0+\overline0\cdot1=0 .
\]

El mismo resultado se obtiene directamente de la definición $\bk fg=\frac1\pi\bigl[(f,g)_{L^2}+i\,(i\cdot f,g)_{L^2}\bigr]$ \eqref{eq:hermitico}, lo que
muestra el papel de cada término. Como $i\cdot\cos nt=\sin nt$, por $(\cos nt,\cos mt)_{L^2}=\pi\delta_{nm},\quad\penalty0 (\sin nt,\sin mt)_{L^2}=\pi\delta_{nm},\quad\penalty0 (\cos nt,\sin mt)_{L^2}=0$ \eqref{eq:ortogonalidad}
se tiene
\[
  \bk00=\frac1\pi\Bigl[\Lp{\cos t}{\cos t}+i\,\Lp{\sin t}{\cos t}\Bigr]
  =\frac1\pi\bigl[\pi+i\cdot0\bigr]=1,
\]
\[
  \bk11=\frac1\pi\Bigl[\Lp{\cos 2t}{\cos 2t}+i\,\Lp{\sin 2t}{\cos 2t}\Bigr]
  =\frac1\pi\bigl[\pi+i\cdot0\bigr]=1,
\]
\[
  \bk01=\frac1\pi\Bigl[\Lp{\cos t}{\cos 2t}+i\,\Lp{\sin t}{\cos 2t}\Bigr]
  =\frac1\pi\bigl[0+i\cdot0\bigr]=0 .
\]
En los dos primeros cálculos se ve que el factor $1/\pi$ es el que lleva la norma de $\ket0$ y
$\ket1$ al valor $1$.

\paragraph{Conclusión.}
Con la suma usual de funciones y la multiplicación por escalares complejos
$z\cdot f=(z\alpha)\cdot\cos t+(z\beta)\cdot\cos 2t$ \eqref{eq:mult}, que extiende la regla del enunciado, $V$ es un espacio vectorial complejo, y
$\{\ket0,\ket1\}=\{\cos t,\cos 2t\}$ es una base, de modo que $\dim_{\mathbb C}V=2$. A partir del
producto interno de $L^2$ se construyó el producto hermítico
\[
  \destacado{\bk fg=\frac1\pi\Bigl[\Lp fg+i\,\Lp{i\cdot f}{g}\Bigr]
  =\overline\alpha\,\alpha'+\overline\beta\,\beta' ,}
\]
respecto al cual $\{\ket0,\ket1\}$ es una base ortonormal. En coordenadas complejas, $V$ con
este producto se comporta exactamente como $\mathbb C^2$ con su producto usual, que es el
espacio de estados de un qubit.
% ======================== FIN DEL PUNTO 1.1 ============================